\chapter{Quỹ Dự Phòng Là Bình Yên}
\markboth{Quỹ Dự Phòng Là Bình Yên}{An Emergency Fund Is Peace}

\section{Một Đêm Con Sốt}
\begin{truyen}{Khi Biến Cố Đến Không Hẹn Trước}{When Trouble Arrives Unannounced}
\chuhoa{M}{ột đêm khuya, đứa con lên cơn sốt cao. Người cha bế con ra xe, đi trong tâm trạng vừa lo vừa vội. Những khoản tiền dành cho bệnh viện, thuốc men, xét nghiệm bỗng hiện ra dồn dập.}
\textit{(Late one night, the child developed a high fever. The father carried the child to the vehicle in a state of fear and urgency. Hospital fees, medicine, and tests suddenly rose before him all at once.)}

Trong những lúc như vậy, người ta mới thấy tiền để dành không phải là thứ xa xỉ. Nó là thứ giúp trái tim đang sợ có chỗ bấu vào. Nếu không có khoản nhỏ cất sẵn, chỉ một cơn bệnh cũng đủ biến thành cơn hoảng loạn tài chính.
\textit{(At such moments, one sees clearly that saved money is not a luxury. It is what gives a frightened heart something to hold on to. Without a small reserve set aside, even one illness can become a financial panic.)}

Người thầy mở phong bì dự phòng đã cất từ trước. Không nhiều, nhưng đủ. Cảm giác lúc ấy không phải sung sướng. Chỉ là biết ơn một quyết định cũ của chính mình.
\textit{(The teacher opened the emergency envelope he had prepared beforehand. It was not much, but it was enough. The feeling at that moment was not joy. It was gratitude toward an earlier wise decision of his own.)}
\end{truyen}

\begin{baihoc}
Quỹ dự phòng không làm ta giàu. Nó giúp ta không ngã gục khi đời bất ngờ đẩy mạnh một cái.
\textit{(An emergency fund does not make us rich. It helps us avoid collapse when life suddenly pushes hard.)}
\end{baihoc}

\begin{ghinhoanh}
Emergency money is stored calm.

It waits silently for difficult days.
\end{ghinhoanh}

\begin{tuhocnhanh}
1. Nếu tối nay có chuyện gấp, em có khoản nào để xoay không? \textit{If something urgent happened tonight, do you have money set aside?}

2. Em đang trì hoãn quỹ dự phòng vì lý do gì? \textit{Why are you postponing an emergency fund?}
\end{tuhocnhanh}
\ngancach

\section{Tích Cốc Phòng Cơ}
\begin{truyen}{Chuẩn Bị Không Phải Là Bi Quan}{Preparation Is Not Pessimism}
\chuhoa{N}{gười xưa dạy ``tích cốc phòng cơ, tích y phòng hàn'' --- tích trữ lúa để phòng đói, dành sẵn áo để phòng rét. Đây không phải lối sống sợ hãi, mà là cách sống có trí nhớ.}
\textit{(The ancients taught, ``Store grain against hunger; store clothing against cold.'' This is not a fearful lifestyle, but a way of living with memory.)}

Ai từng thấy gia đình khổ vì một cơn đau, một vụ mùa thất, một việc làm gián đoạn, người ấy sẽ hiểu chuẩn bị trước là một dạng thương người rất cụ thể. Không chờ tới lúc có chuyện mới hoảng hốt đi mượn vay khắp nơi.
\textit{(Whoever has seen a family suffer from illness, crop failure, or interrupted work will understand that preparation is a very concrete form of love. It means not waiting until trouble comes to borrow in panic from all directions.)}

Sống ngày nào biết dành một ít cho ngày mai là đang tôn trọng sự vô thường của đời sống. Thừa nhận rằng biến cố có thể đến không làm ta yếu đi. Nó làm ta tỉnh táo hơn.
\textit{(To set aside a little each day for tomorrow is to respect life’s uncertainty. Admitting that hardship may come does not make us weak. It makes us clearer.)}
\end{truyen}

\begin{baihoc}
Người chuẩn bị trước không phải vì bi quan, mà vì yêu gia đình mình đủ nhiều để không phó mặc.
\textit{(Those who prepare are not pessimistic; they simply love their families too much to leave them unprotected.)}
\end{baihoc}

\begin{ghinhoanh}
Preparation is practical love.

It is not fear.
\end{ghinhoanh}

\begin{tuhocnhanh}
1. Chuẩn bị trước điều gì là biểu hiện của yêu thương? \textit{How is preparation an act of love?}

2. Em đang phó mặc điều gì cho may rủi? \textit{What are you currently leaving to chance?}
\end{tuhocnhanh}
\ngancach

\section{Phong Bì Đứng Gác Cho Gia Đình}
\begin{truyen}{Một Khoản Tiền Nằm Yên Nhưng Không Vô Ích}{Money That Lies Still but Is Not Useless}
\chuhoa{N}{hiều người nhìn tiền để dành mà sốt ruột. Họ nghĩ số tiền ấy nằm im, không sinh lời, như thế là lãng phí cơ hội.}
\textit{(Many people grow restless when they look at saved money. They think the money is lying still, earning little, and that this means wasted opportunity.)}

Nhưng người thầy hiểu phong bì dự phòng không sinh ra để chạy đua. Nó sinh ra để đứng gác. Vai trò của nó không phải làm người ta thấy mình giỏi đầu tư, mà để gia đình có chỗ dựa khi bất an kéo đến.
\textit{(But the teacher understood that the emergency envelope was not born to race. It was born to stand guard. Its role was not to make him feel like a smart investor, but to give the family something to lean on when fear arrived.)}

Trong đời sống bình thường, có những đồng tiền phải chạy để sinh ra thêm giá trị. Nhưng cũng có những đồng tiền phải đứng yên để giữ cánh cửa khỏi bật tung. Biết phân vai cho tiền là một dạng trưởng thành.
\textit{(In ordinary life, some money should move in order to create more value. But some money must stand still to keep the door from flying open. Knowing how to assign roles to money is a form of maturity.)}
\end{truyen}

\begin{baihoc}
Không phải khoản tiền nào cũng cần sinh lời ngay. Có khoản được sinh ra để giữ bình an.
\textit{(Not every amount of money must generate profit immediately. Some amounts exist to preserve peace.)}
\end{baihoc}

\begin{ghinhoanh}
Some money must grow.

Some money must guard.
\end{ghinhoanh}

\begin{tuhocnhanh}
1. Em có đang bắt mọi khoản tiền phải ``làm việc'' mà quên vai trò bảo vệ không? \textit{Are you forcing every coin to work while forgetting its protective role?}

2. Khoản nào trong đời em cần đứng gác từ hôm nay? \textit{What fund in your life needs to start guarding today?}
\end{tuhocnhanh}

\begin{turanminh}
1. Tôi không được lấy lý do thu nhập chưa cao để trì hoãn quỹ dự phòng mãi mãi.

2. Mỗi tháng còn giữ lại được một phần nhỏ là tôi đang mua trước sự bình tĩnh cho những ngày bất ngờ.

3. Người cha không chỉ lo hôm nay đủ ăn; người cha còn phải lo khi con ốm, khi nhà có biến cố, mình không bị cuống lên vì không còn gì trong tay.
\end{turanminh}

\begin{dayhoc}
1. Với học sinh, có thể chuyển bài học ``quỹ dự phòng'' thành khái niệm dễ hiểu hơn: phải biết để dành, không tiêu sạch mọi thứ vừa có.

2. Thầy có thể kể cho các em hiểu rằng trong gia đình, một khoản nhỏ để dành đôi khi cứu được cả một cơn hoảng loạn.

3. Câu hỏi tốt cho lớp là: ``Vì sao người biết chuẩn bị trước thường bình tĩnh hơn khi gặp chuyện?''
\end{dayhoc}